\chapter{Shannon / Perfect Secrecy}

\section{Definitions}
Last chapter, we introduced the concept and reason for an SKE.

\begin{definition}[Symmetric Key Encryption (SKE)]
  A \textbf{symmetric key encryption scheme} is a tuple of algorithms $(\Gen, \Enc, \Dec)$ such that:
  \begin{itemize}
    \item $\Gen$ is a probabilistic algorithm that outputs a key $k$ from the key space $\mathcal{K}$.
    \item $\Enc$ is a probabilistic algorithm that takes a key $k$ and a message $m$ from the message space $\mathcal{M}$ and outputs a ciphertext $c$ from the ciphertext space $\mathcal{C}$.
    \item $\Dec$ is a deterministic algorithm that takes a key $k$ and a ciphertext $c$ and outputs a message $m$ or an error symbol $\perp$.
  \end{itemize}
  The scheme must satisfy the correctness property: for all keys $k \in \mathcal{K}$ and messages $m \in \mathcal{M}$, we have $\Dec(k, \Enc(k, m)) = m$.
\end{definition}

For security, there're a lot of ways to define it. We might say the scheme is safe if the adversary cannot recover 90\% of the message, or if the adversary can't distinguish $c$ from random noise.
Here we adopt Shannon's idea:

\begin{remark}
  The secrecy of a message means that the ciphertext should reveal nothing about the message that the adversary didn't already know.
\end{remark}

In other words, the ``posterior'' distribution of the message (conditioned on the ciphertext) should be the same as the ``prior'' distribution of the message. This leads to the following definition:

\begin{definition}[Shannon Secrecy]
  For any probability distribution $D$ over $\mathcal{M}$, an SKE scheme $(\Gen, \Enc, \Dec)$ is said to achieve \textbf{Shannon secrecy} if for every message $\overline{m} \in \mathcal{M}$ and every ciphertext $\overline{c}$ in the range of $\Enc(k, m)$ for some $k \gets \Gen$ and $m \gets D$, we have
  \[
    \Pr_{m\gets D, k \gets \Gen}[m = \overline{m} \mid \Enc(k,m) = \overline{c}] = \Pr_{m \gets D}[m = \overline{m}]
  \]
\end{definition}

Now we play with this definition:

\begin{align*}
  % \LHS &= \Pr_{m\gets D, k \gets \Gen}[m = \overline{m} \mid \Enc(k,m) = \overline{c}] \\
    \LHS   &= \frac{\Pr_{m, k}[m = \overline{m} \wedge \Enc(k,\overline{m}) = \overline{c}]}{\Pr_{m, k}[\Enc(k,m) = \overline{c}]} \\
       &= \frac{\Pr_{m}[m = \overline{m}] \cdot \Pr_{k}[\Enc(k,\overline{m}) = \overline{c}]}{\Pr_{m,k}[\Enc(k,m) = \overline{c}]} \quad \text{(by independence of $k$)} \\
  % \RHS &= \Pr_{m}[m = \overline{m}]
\end{align*}
To let $\LHS = \RHS$, we need
\[ \Pr_{k\gets \Gen}[\Enc(k,\overline{m}) = \overline{c}] = \Pr_{m\gets D,k\gets Gen}[\Enc(k,m) = \overline{c}] \]
For a same distribution $D$, the RHS is a constant. So we can say that for any $\overline{m}$ and $\overline{c}$, the probability of $\overline{c}$ being the ciphertext of $\overline{m}$ is a constant. This means that for any two messages $m_0$ and $m_1$, the probability of $\overline{c}$ being the ciphertext of $m_0$ is equal to that of $m_1$. This leads to a second definition:

\begin{definition}[Perfect Secrecy]
  An SKE scheme $(\Gen, \Enc, \Dec)$ is said to achieve \textbf{perfect secrecy} if for every pair of messages $m_0, m_1 \in \mathcal{M}$ and every ciphertext $\overline{c}\in \mathcal C$, we have
  \[
    \Pr_{k\gets \Gen}[\Enc(k,m_0) = \overline{c}] = \Pr_{k\gets \Gen}[\Enc(k,m_1) = \overline{c}]
  \]
\end{definition}

In other words, the ciphertext has a same distribution regardless of the message.

\begin{theorem}
  Shannon secrecy is equivalent to perfect secrecy.
\end{theorem}
Here we have already proved the $\implies$ direction. The other direction is left as an exercise.

\section{One-time Pad}

The Shannon / perfect secrecy is achievable.

\begin{construction}[One-time Pad (OTP) / Vernam Cipher]
  Let $\mathcal{M} = \mathcal{K} = \mathcal{C} = \{0,1\}^n$ for some $n$. The one-time pad scheme is defined as follows:
  \begin{itemize}
    \item $\Gen$: Choose a key $k$ uniformly at random from $\{0,1\}^n$.
    \item $\Enc(k, m)$: Compute the ciphertext $c = m \oplus k$, where $\oplus$ denotes bitwise XOR.
    \item $\Dec(k, c)$: Compute the message $m = c \oplus k$.
  \end{itemize}
\end{construction}

Correctness holds as $\forall k, m: \Dec(k, \Enc(k, m)) = (m \oplus k) \oplus k = m$.

\begin{theorem}
  The OTP scheme achieves perfect secrecy.
\end{theorem}
\begin{proof}
  Fix $m_0, m_1, \overline{c} \in \{0,1\}^n$. We have
  \[
    \Pr_{k\gets \Gen}[\Enc(k,m_0) = \overline{c}] = \Pr_{k\gets \Gen}[m_0 \oplus k = \overline{c}] = \Pr_{k\gets \Gen}[k = m_0 \oplus \overline{c}] = 2^{-n}
  \]
  Symmetrically, the same probability holds for $m_1$. Thus, the OTP scheme achieves perfect secrecy.
\end{proof}

We even have a stronger result: regardless of the message distribution, the ciphertext is uniformly random over the choice of $k$.

\section{The Downsides}

OTP has some downsides:
\begin{itemize}
  \item $\mathcal K = \mathcal M$, which means the key length must be at least as long as the message length. This is impractical for long messages.
  \item The key can only be used once. Reusing the key for multiple messages compromises security, as 
    \[ \Enc(k, m_0) \oplus \Enc(k, m_1) = (m_0 \oplus k) \oplus (m_1 \oplus k) = m_0 \oplus m_1 \]
    which reveals information about the messages.
  \end{itemize}

  \begin{theorem}
    In any Shannon / perfect secrecy scheme,
    \[
      |\mathcal K| \geq |\mathcal M|
    \]
  \end{theorem}
  \begin{proof}
    Suppose $|\mathcal K| < |\mathcal M|$.

    Take $D$ to be uniform over $\mathcal M$. Suppose we get a ciphertext $\overline{c}$ in the range of $\Enc$.

    Consider $P = \{\Dec(k,\overline{c}) : k \in \mathcal K\}$. Any message that can produce $\overline{c}$ must be in $P$, by correctness.

    However, as $|P| \leq |\mathcal K|$ (as each key produces at most one message), we have $|P| < |\mathcal M|$. Thus, there exists a message $m^* \in \mathcal M\setminus P$.

    This means that $\Pr_{m\gets D, k\gets \Gen}[m = m^* \mid \Enc(k,m) = \overline{c}] = 0$, while $\Pr_{m\gets D}[m = m^*] = 1/|\mathcal M| > 0$. This contradicts the definition of Shannon secrecy. Therefore, we must have $|\mathcal K| \geq |\mathcal M|$.
  \end{proof}
